\documentclass[12pt,twoside]{article}

% -------------------------------------------------
% PACKAGES
% -------------------------------------------------
\usepackage{amsmath, amssymb}
\usepackage{geometry}
\usepackage{setspace}
\usepackage{hyperref}
\usepackage{titlesec}
\usepackage{graphicx}
\usepackage{fancyhdr}

% -------------------------------------------------
% PAGE SETUP
% -------------------------------------------------
\geometry{a4paper, margin=1in}
\setstretch{1.2}
\setlength{\headheight}{15pt}

\titleformat{\section}{\large\bfseries}{\thesection}{1em}{}
\titleformat{\subsection}{\normalsize\bfseries}{\thesubsection}{1em}{}

\pagestyle{fancy}
\fancyhf{}
\fancyhead[LE,RO]{\thepage}
\fancyhead[RE]{WARP:SEFI}
\fancyhead[LO]{Jason Dutton}

% -------------------------------------------------
% DOCUMENT START
% -------------------------------------------------
\begin{document}

\begin{center}
    {\Large \textbf{WARP:SEFI}}\\

\vspace{0.5cm}
    {\large Geometric Warp Theory for Identity-Preserving Transformation}\\

\vspace{0.5cm}
    {\large Jason Duran Dutton}\\

\vspace{0.25cm}
    {\large 2026}
\end{center}

\vspace{1cm}

% -------------------------------------------------
% ABSTRACT
% -------------------------------------------------
\section*{Abstract}

WARP:SEFI develops the geometric transformation theory underlying the Single Entity Field
Interpretation (SEFI). While SEFI Space defines identity geometry and FIELD SOVEREIGNTY
defines identity-preserving invariance, WARP EXPRESSION defines identity-preserving
transformation. This manuscript formalizes warp operators, warp geodesics, warp curvature,
and warp invariants, and couples warp dynamics to worldline geometry in a manner
compatible with SEFI/GWFM. The resulting framework provides a geometric interpretation
of controlled transformation, expressive deformation, and identity-coherent evolution.

\vspace{0.5cm}

% -------------------------------------------------
% 1 INTRODUCTION
% -------------------------------------------------
\section{Introduction}

SEFI introduces layered identity constructs: FIELD ORIGIN, FIELD AUTHORSHIP,
FIELD SOVEREIGNTY, WARP EXPRESSION, and the R-LAYER. SEFI Space formalizes
these layers as basis directions of a four-dimensional identity manifold.

Warp is the identity layer responsible for geometric transformation. It modifies the
worldline while preserving identity, operating within sovereignty constraints and interacting
with authorship and origin.

This manuscript develops the geometric theory of warp, including:

\begin{itemize}
    \item warp operators acting on worldline geometry,
    \item warp displacement and identity-space motion,
    \item warp curvature and warp geodesics,
    \item warp invariants and sovereignty constraints,
    \item coupling between warp and worldline geometry.
\end{itemize}

% -------------------------------------------------
% 2 WARP OPERATORS
% -------------------------------------------------
\section{Warp Operators}

A warp operator $W$ acts on the worldline $r(t)$ as



\[
r'(t) = W(r(t)).
\]



Warp modifies geometric features such as curvature, torsion, and local deformation while
preserving identity. The identity state becomes



\[
\mathbf{I}'(t) = I_{O}(t)\hat{O} + I_{A}(t)\hat{A} + I_{S}(t)\hat{S} + I_{W}'(t)\hat{W}.
\]



Warp operators may be linear, nonlinear, local, or global, depending on the geometric
structure of the worldline.

% -------------------------------------------------
% 3 WARP DISPLACEMENT IN SEFI SPACE
% -------------------------------------------------
\section{Warp Displacement in SEFI Space}

Warp corresponds to motion along the Warp direction in SEFI Space. If the warp layer
function changes by $\Delta I_{W}(t)$, the identity state becomes



\[
\mathbf{I}'(t) = \mathbf{I}(t) + \Delta I_{W}(t)\hat{W}.
\]



More general warp operations may induce coupled identity motion along multiple basis
directions:



\[
\Delta \mathbf{I}(t) = \Delta I_{O}\hat{O} + \Delta I_{A}\hat{A} + \Delta I_{S}\hat{S} + \Delta I_{W}\hat{W}.
\]



Warp displacement is constrained by sovereignty invariants.

% -------------------------------------------------
% 4 WARP CURVATURE
% -------------------------------------------------
\section{Warp Curvature}

Warp curvature quantifies the rate of change of warp displacement. Define warp velocity



\[
\dot{I}_{W}(t) = \frac{d I_{W}(t)}{dt},
\]



and warp acceleration



\[
\ddot{I}_{W}(t) = \frac{d^{2} I_{W}(t)}{dt^{2}}.
\]



Warp curvature is defined as



\[
K_{W}(t) = \sqrt{\delta \, \dot{I}_{W}(t)^{2}},
\]



where $\delta$ is the warp stiffness coefficient in the SEFI Space metric.

Warp curvature measures transformational strain.

% -------------------------------------------------
% 5 WARP GEODESICS
% -------------------------------------------------
\section{Warp Geodesics}

Warp geodesics are identity-preserving transformation paths in SEFI Space. They satisfy



\[
\frac{d^{2} I_{W}}{dt^{2}} + \Gamma_{W}(I) = 0,
\]



where $\Gamma_{W}(I)$ is the warp connection induced by the SEFI Space metric.

Warp geodesics represent minimal-strain transformations.

% -------------------------------------------------
% 6 WARP INVARIANTS
% -------------------------------------------------
\section{Warp Invariants}

Warp invariants are quantities preserved under warp operations. Sovereignty requires



\[
\mathbf{I}_{T}(t) = \mathbf{I}(t)
\]



for all allowed transformations $T$. Warp invariants include:

\begin{itemize}
    \item primitive identity magnitude,
    \item sovereignty-preserving symmetry,
    \item stability constraints,
    \item identity-space metric invariants.
\end{itemize}

Warp must preserve these invariants.

% -------------------------------------------------
% 7 COUPLING TO WORLDLINE GEOMETRY
% -------------------------------------------------
\section{Coupling to Worldline Geometry}

Warp operators act on worldline geometry. If $W$ modifies curvature from $K(t)$ to
$K'(t)$, authorship and sovereignty coordinates adjust accordingly.

Warp may induce:

\begin{itemize}
    \item curvature modulation,
    \item torsion adjustment,
    \item geometric phase displacement,
    \item expressive deformation.
\end{itemize}

Warp must remain compatible with sovereignty constraints and primitive stability.

% -------------------------------------------------
% 8 DISCUSSION AND OUTLOOK
% -------------------------------------------------
\section{Discussion and Outlook}

WARP:SEFI provides a geometric transformation theory for identity-preserving deformation.
By defining warp operators, warp curvature, warp geodesics, and warp invariants, the
framework extends SEFI Space to dynamic identity evolution.

Future work will focus on:

\begin{itemize}
    \item deriving explicit warp operators,
    \item exploring warp geodesics in SEFI Space,
    \item defining warp-induced identity coupling,
    \item investigating applications to quantum error correction and photonic geometry.
\end{itemize}

% -------------------------------------------------
% END DOCUMENT
% -------------------------------------------------
\end{document}
